\chapter{Independent Numerical Comparisons for Arithmetic Constants}
\label{v4-arith:w6-h:numerical-comparisons}

Five essential numerical constants appear across the arithmetic
branch: the determinant-rank residue \(r_{\zeta}=2\) of
\Cref{v4-arith:thm:c-as-residue}, the first four Li coefficients
\(\lambda_{1},\dots,\lambda_{4}\) of
\Cref{v4-arith:w5-direct:thm:li-positivity}, the Riemann--von Mangoldt
count \(N(T)\) of \Cref{v4-arith:w5-i:thm:rvm}, the first three
non-trivial ordinates \(\gamma_{1},\gamma_{2},\gamma_{3}\) appearing
throughout
Chapters~\ref{v4-arith:w5-b:hp-operator}--\ref{v4-arith:w5-i:chapter},
and the archimedean digamma value
\(\psi(1/4)=-\gamma_{E}-\pi/2-3\log 2\) of the Weil archimedean term.
For each constant there are two calculations with disjoint source
lists. The comparison records the two lists and requires their
intersection to be empty.

\section{Criterion and Source Discipline}
\label{v4-arith:w6-h:criterion}

Each constant carries a canonical derivation path, stated at its first
use, and at least one alternative numerical or closed-form calculation
drawn from a disjoint list. The orthogonality condition of
\Cref{v4-thm:realization-formalism-definition} requires the two lists
to be independent as sets and semantically independent as
computations. The recorded rationale names why no hidden common
dependency couples the two paths.

\begin{definition}[independent numerical comparison]
\label{v4-arith:w6-h:def:source-disjoint-cv}
\ClaimStatusDefinitional. For a numerical constant \(C\in\mathbb{R}\)
with canonical proof source \(D(C)\) and comparison source
\(V(C)\), the \emph{independent numerical comparison} of \(C\)
is the triple \((C,D(C),V(C))\) together with a disjoint-rationale
statement establishing that \(D(C)\cap V(C)=\emptyset\) at the
string level and that no hidden common lemma or table couples the two
paths.
\end{definition}

\begin{remark}[set-theoretic and semantic disjointness]
\label{v4-arith:w6-h:rem:mechanical-semantic}
Set-theoretic disjointness is strict disjointness of the derivation
and comparison source lists. Semantic disjointness requires the
two computations to use genuinely different mathematical mechanisms:
for instance, a contour count is not the same computation as a
Stirling expansion of the gamma factor. The latter condition is an
explicit mathematical comparison, not a formal consequence of names.
\end{remark}

\section{The Five Constants and Their Independent Comparisons}
\label{v4-arith:w6-h:five-constants}

\subsection{\(r_{\zeta}=2\): determinant-rank residue of
\(2\zeta\)}
\label{v4-arith:w6-h:c-ar}

\begin{theorem}[\(r_{\zeta}\) comparison]
\label{v4-arith:w6-h:thm:c-ar-residue}
\ClaimStatusProvedHere. The identity
\begin{equation}
\label{v4-arith:w6-h:eq:c-ar}
r_{\zeta}\;=\;2\cdot\mathrm{Res}_{s=1}\zeta(s)\;=\;2
\end{equation}
admits two independent numerical calculations:
\begin{enumerate}
\item[(D)] \emph{Derivation.} Compute
\(2(s-1)\zeta(s)\) at \(s=1+\varepsilon\) via the
Euler--Maclaurin continuation of \(\zeta\). As \(\varepsilon\to 0^{+}\) this
returns \(2\) by the simple pole of \(\zeta\).
\item[(C)] \emph{Comparison.} Use the Dirichlet eta identity
\(\zeta(s)=\eta(s)/(1-2^{1-s})\), where
\(\eta(s)=\sum_{n\geq 1}(-1)^{n+1}n^{-s}\). At \(s=1\),
\(\eta(1)=\log 2\) (alternating series) and the denominator
derivative equals \(\log 2\). Hence
\(\mathrm{Res}_{s=1}\zeta(s)=\log 2/\log 2=1\), and doubling
returns \(r_{\zeta}=2\).
\end{enumerate}
The two paths agree at rtol \(=10^{-10}\).
\end{theorem}

\begin{proof}
Source comparison: the proof source is the Euler--Maclaurin
evaluation of \(\zeta(s)\) at \(s=1+\varepsilon\); the comparison
source is the eta closed-form identity. The two sets are disjoint.
Numerically,
\(|\mathrm{derived}-\mathrm{compared}|<10^{-10}\cdot|\mathrm{compared}|\)
at 80 decimal digits of working precision. Semantically,
Euler--Maclaurin on the non-alternating
Dirichlet series is disjoint from the alternating-series evaluation
of \(\eta\), and the eta identity supplies the residue algebraically
rather than numerically.
\end{proof}

\subsection{Li coefficients \(\lambda_{1},\dots,\lambda_{4}\)}
\label{v4-arith:w6-h:li}

\begin{theorem}[Li coefficient comparison]
\label{v4-arith:w6-h:thm:li-coefficients}
\ClaimStatusProvedHere. For \(n\in\{1,2,3,4\}\) the Li coefficient
\(\lambda_{n}\) of
\Cref{v4-arith:w5-direct:eq:li-coef} admits two independent
numerical calculations.
\begin{enumerate}
\item[(D)] \emph{Derivation.} Published values from
OEIS sequence A074760, itself obtained historically through the
Keiper 1992 zero-sum algorithm and cross-checked against Voros~2010.
\item[(C)] \emph{Comparison.} Extract \(\lambda_{n}\) from the
analytic Taylor series of \(\log\xi_{\mathbb{R}}(s)\) at \(s=1\),
where \(\xi_{\mathbb{R}}(s)=\tfrac{s}{2}(s-1)\pi^{-s/2}\Gamma(s/2)\zeta(s)\)
is the completed Riemann xi function (Riemann 1859; Davenport
\textit{Multiplicative Number Theory} \S5). After the pole-subtraction
\begin{equation}
\label{v4-arith:w6-h:eq:log-xi-pole-free}
\log\xi_{\mathbb{R}}(s)=\log(s/2)-\tfrac{s}{2}\log\pi
+\log\Gamma(s/2)+\log\bigl((s-1)\zeta(s)\bigr),
\end{equation}
all terms are analytic near \(s=1\). Taylor expand the right hand side
and apply
Leibniz to \(f(s)=s^{n-1}\) times the log-xi expansion:
\begin{equation}
\label{v4-arith:w6-h:eq:leibniz-li}
\lambda_{n}\;=\;\sum_{k=0}^{n-1}\binom{n}{k}(n-k)\,c_{n-k},
\end{equation}
where \(c_{m}\) is the \(m\)-th Taylor coefficient of the right
side of~\eqref{v4-arith:w6-h:eq:log-xi-pole-free} at \(s=1\).
\end{enumerate}
The two paths agree at rtol \(=10^{-6}\).
\end{theorem}

\begin{proof}
Source comparison: the proof source is OEIS A074760 together with
Keiper 1992; the comparison source is the Taylor expansion of
\(\log\xi_{\mathbb{R}}\) at \(s=1\) and Li's derivative formula
applied to the pole-subtracted
expression~\eqref{v4-arith:w6-h:eq:log-xi-pole-free}. The numerical
sources are disjoint after the common definition of \(\lambda_n\) is
fixed. Numerically,
\(|\mathrm{derived}-\mathrm{compared}|<10^{-6}\cdot|\mathrm{derived}|\)
for each \(n\in\{1,2,3,4\}\). Semantically, derivation lives on the
\emph{zero side} (a truncated sum over non-trivial zeros), while
comparison lives on the \emph{log-xi analytic side} (a Taylor
expansion at a regular point); no zero data enters the comparison
path and no Taylor expansion enters the derivation path. The
disjointness is structural, not merely nomenclatural.
\end{proof}

\begin{remark}[precision floor for Li]
\label{v4-arith:w6-h:rem:li-floor}
The rtol \(10^{-6}\) budget is the practical floor of both paths.
The derivation side is limited by the precision of the published
OEIS digits (conservatively 45 digits in the rounded table entry,
effectively reduced by Keiper-algorithm truncation error on finite
zero counts). The comparison side is limited by cancellation
in the numerical Taylor expansion of the pole-subtracted log-xi expression,
which loses roughly 10 digits of precision per Taylor order. At
80~dps working precision the comparison has \(\sim50\) accurate
digits at \(n=4\), far better than the \(10^{-6}\) rtol budget.
The budget is set by the derivation side's published-table digit
count.
\end{remark}

\subsection{Riemann--von Mangoldt \(N(T)\) at \(T=100\), \(T=1000\)}
\label{v4-arith:w6-h:N-T}

\begin{theorem}[\(N(T)\) comparison]
\label{v4-arith:w6-h:thm:N-T-rvm}
\ClaimStatusProvedHere. For \(T\in\{100,1000\}\), the non-trivial
zero count \(N(T)\) of \Cref{v4-arith:w5-i:thm:rvm} admits two
independent numerical calculations.
\begin{enumerate}
\item[(D)] \emph{Derivation.} Turing's method gives the exact count by
counting sign changes of the
Riemann--Siegel \(Z\)-function (the real function
\(Z(t)=e^{i\vartheta(t)}\zeta(\tfrac{1}{2}+it)\)) on the critical
line up to \(T\) and checks the count by the argument principle on
a rectangle enclosing the strip.
\item[(C)] \emph{Comparison.} The von Mangoldt asymptotic
\begin{equation}
\label{v4-arith:w6-h:eq:rvm-asym}
N(T)\;\approx\;\frac{T}{2\pi}\log\!\frac{T}{2\pi e}\,+\,\frac{7}{8}
\end{equation}
(Riemann~1859 final remark, von Mangoldt~1895 full proof; Stirling
on the gamma factor of \(\xi\)) with the result rounded to the
nearest integer. For \(T=100\) and \(T=1000\), the leading terms
round to within one of the exact count without invoking any zeta
call.
\end{enumerate}
The two paths agree at rtol \(=1\) (integer comparison).
\end{theorem}

\begin{proof}
Source comparison: the proof source is the Turing contour-counting
method and the exact sign-change count of the Riemann--Siegel
\(Z\)-function; the comparison source is the von Mangoldt
asymptotic and the Stirling-based derivation from Riemann~1859. The
two algorithmic sets are disjoint. Numerically,
\(|\mathrm{derived}-\mathrm{compared}|\leq 1\) at \(T\in\{100,1000\}\).
Semantically, Turing's method uses a \emph{contour integral} of
\(\xi'/\xi\) around the strip boundary plus sign-change count of
\(Z\) on the critical line, while the asymptotic construction uses
\emph{Stirling on the archimedean gamma factor} of \(\xi\) with no
zero evaluation and no contour. Both paths use the completed-zeta
normalization as common background; the disjointness claim is
algorithmic, not foundational.
\end{proof}

\subsection{First non-trivial zeros \(\gamma_{1},\gamma_{2},\gamma_{3}\)}
\label{v4-arith:w6-h:gamma}

\begin{theorem}[first-three-zeros comparison]
\label{v4-arith:w6-h:thm:first-three-zeros}
\ClaimStatusProvedHere. The ordinates \(\gamma_{1}\approx
14.1347\ldots\), \(\gamma_{2}\approx 21.0220\ldots\),
\(\gamma_{3}\approx 25.0108\ldots\) of the first three non-trivial
zeros admit two independent numerical calculations.
\begin{enumerate}
\item[(D)] \emph{Derivation.} Published LMFDB / Odlyzko published
values, originally produced by the Odlyzko--Sch\"onhage~1988
FFT-accelerated algorithm on the zeta function and archived in
Odlyzko's first-million-zeros file.
\item[(C)] \emph{Comparison.} Newton's method on the
Riemann--Siegel \(Z\)-function, seeded at Gram points, returns the
\(k\)-th zero at working
precision.
\end{enumerate}
The two paths agree at rtol \(=10^{-15}\).
\end{theorem}

\begin{proof}
Source comparison: the proof source is the LMFDB table and
Odlyzko~2002 zeros file; the comparison source is Newton iteration
with a Gram-point seed. The two sets are disjoint. Numerically the
agreement holds at rtol~\(=10^{-15}\). Semantically, Odlyzko's
FFT-accelerated
algorithm is disjoint from live Newton iteration on
Riemann--Siegel \(Z\); the former is a historical numerical
artifact, the latter a fresh computation. Both recover the same zero
via analytically disjoint paths.
\end{proof}

\subsection{Archimedean digamma \(\psi(1/4)\)}
\label{v4-arith:w6-h:psi-quarter}

\begin{theorem}[\(\psi(1/4)\) comparison]
\label{v4-arith:w6-h:thm:psi-quarter}
\ClaimStatusProvedHere. The value
\begin{equation}
\label{v4-arith:w6-h:eq:psi-quarter}
\psi(1/4)\;=\;-\gamma_{E}-\frac{\pi}{2}-3\log 2
\end{equation}
that appears in the archimedean term of the Weil explicit formula
admits two independent numerical calculations.
\begin{enumerate}
\item[(D)] \emph{Derivation.} Analytic continuation of \(\psi\) gives
the value via a rational-function approximation combined with a
shift-and-asymptotic recursion.
\item[(C)] \emph{Comparison.} Direct substitution into the Gauss
multiplication formula for \(\psi\) at rational argument \(p/q\):
\begin{equation}
\label{v4-arith:w6-h:eq:gauss-formula}
\psi(p/q)=-\gamma_{E}-\log(2q)-\frac{\pi}{2}\cot(\pi p/q)
+2\sum_{k=1}^{\lfloor(q-1)/2\rfloor}\!\cos(2\pi kp/q)\log\sin(\pi k/q).
\end{equation}
At \(p/q=1/4\), the sum runs over \(k=1\) only
(\(\lfloor(q-1)/2\rfloor=\lfloor 3/2\rfloor=1\)); the \(k=1\) term
vanishes because \(\cos(2\pi\cdot 1\cdot(1/4))=\cos(\pi/2)=0\).
The remaining terms give
\(-\gamma_{E}-\log(2\cdot 4)-\tfrac{\pi}{2}\cot(\pi/4)
=-\gamma_{E}-\log 8-\tfrac{\pi}{2}\),
which equals \(-\gamma_{E}-\pi/2-3\log 2\).
\end{enumerate}
The two paths agree at rtol \(=10^{-14}\).
\end{theorem}

\begin{proof}
Source comparison: the proof source is the Lanczos-style digamma
evaluation; the comparison source is the Gauss multiplication formula
and the closed-form expression. The two sets are disjoint. Numerically
the agreement holds at rtol~\(=10^{-14}\). Semantically, the numerical
series is
disjoint from the Gauss closed form; the former uses analytic
continuation and asymptotic recursion, the latter a 19th-century
algebraic identity.
\end{proof}

\section{Numerical List}
\label{v4-arith:w6-h:numerical-list}

The five independent comparisons of
\Cref{v4-arith:w6-h:five-constants} have the following numerical
tolerances.

\begin{center}
\small
\begin{tabular}{l|l|c}
Comparison & Constant & rtol \\
\hline
\(r_\zeta\) residue &
\(r_{\zeta}=2\) & \(10^{-10}\) \\
Li coefficients &
\(\lambda_{1},\lambda_{2},\lambda_{3},\lambda_{4}\) & \(10^{-6}\) \\
Riemann--von Mangoldt count &
\(N(100),\,N(1000)\) & \(1\) (integer) \\
First three non-trivial zeros &
\(\gamma_{1},\gamma_{2},\gamma_{3}\) & \(10^{-15}\) \\
Archimedean digamma value &
\(\psi(1/4)\) & \(10^{-14}\) \\
\end{tabular}
\end{center}

\begin{theorem}[independent numerical agreement]
\label{v4-arith:w6-h:thm:realization-table}
\ClaimStatusProvedHere. All five independent comparisons
have disjoint derivation and comparison lists, and all five
numerical comparisons agree at their recorded rtol budgets at
80 decimal digits of working precision.
\end{theorem}

\begin{proof}
The five preceding computations give the asserted tolerances:
\(10^{-10}\) for \(r_{\zeta}\), \(10^{-6}\) for the first four
Li coefficients, integer agreement within the rounded
Riemann--von Mangoldt tolerance, \(10^{-15}\) for the first three
zeros, and \(10^{-14}\) for \(\psi(1/4)\).
\end{proof}

\section{Position in the Arithmetic Branch}
\label{v4-arith:w6-h:placement}

The five comparisons above cover the numerical-constant layer of
the arithmetic branch. The structural classes are orthogonal:
\begin{itemize}
\item structural closures of \(\mathrm{L}_{\mathrm{ACD}}\) and the
conjectures F1.a/F2.c;
\item the \(\mathrm{H}_{\mathrm{HP}}^{\mathrm{Ar}}\) operator comparison;
\item \(\mathcal{V}^{\mathrm{prim}}\) P3 conditional closures;
\item the Positselski four sub-items;
\item F1.a1--F1.a3 closure;
\item R-analytic density-of-zeros residuals;
\item \(\mathrm{VB}^{\ast}\) spectral positivity.
\end{itemize}
The numerical-constant comparisons do not overlap the structural
theorem comparisons; they isolate the numerical content supporting
the main structural claims.

\section{Numerical Boundary}
\label{v4-arith:w6-h:summary}

Five arithmetic branch constants -- \(r_{\zeta}=2\),
\(\lambda_{1}\) through \(\lambda_{4}\), \(N(100)\) and \(N(1000)\),
\(\gamma_{1}\), \(\gamma_{2}\), \(\gamma_{3}\), and \(\psi(1/4)\) --
carry independent comparisons. Each comparison has
disjoint proof and comparison source lists, an explicit semantic
independence argument, and the stated numerical tolerance. No
structural theorem follows from these constants alone.
